% Generated draft: requires mathematical and source review.
\documentclass[11pt]{article}
\usepackage[margin=25mm]{geometry}
\usepackage{fontspec}
\setmainfont{Latin Modern Roman}
\usepackage{amsmath,amssymb,mathtools,unicode-math}
\setmathfont{Latin Modern Math}
\usepackage{enumitem,needspace,xurl,hyperref}
\setlength{\parindent}{0pt}
\setlength{\parskip}{5pt}
\setlength{\emergencystretch}{3em}
\newcommand{\N}{\mathbb{N}}
\newcommand{\Z}{\mathbb{Z}}
\newcommand{\Q}{\mathbb{Q}}
\newcommand{\R}{\mathbb{R}}
\newcommand{\C}{\mathbb{C}}
\newcommand{\F}{\mathbb{F}}
\newcommand{\A}{\mathbb{A}}
\newcommand{\PP}{\mathbb{P}}
\newcommand{\E}{\mathbb{E}}
\newcommand{\eps}{\varepsilon}
\newcommand{\dd}{\,\mathrm{d}}
\begin{document}
\section*{UnsolvedMath ID: 701630}
\textbf{Source problem code:} catalogue entry
\subsection{Definitions and mathematical statement}
For a finite group \(G\), let \(F(G)\) be its largest nilpotent normal subgroup. A component of \(G\) is a subnormal quasisimple subgroup, where a group \(L\) is quasisimple if \(L=[L,L]\) and \(L/Z(L)\) is nonabelian simple. Here subnormal means connected to \(G\) by a finite chain of normal inclusions. Let \(E(G)\) be the subgroup generated by all components of \(G\). The generalized Fitting subgroup is
\[
F^*(G)=F(G)E(G).
\]
Define \(F^*_0(G)=1\), and, for each integer \(i\geq 0\), let \(\pi_i:G\to G/F^*_i(G)\) be the quotient map and set
\[
F^*_{i+1}(G)=\pi_i^{-1}\bigl(F^*(G/F^*_i(G))\bigr).
\]
The generalized Fitting height is
\[
h^*(G)=\min\{k\in\mathbb Z_{\geq 0}:F^*_k(G)=G\}.
\]
A group \(K\) is soluble if its derived series, defined by \(K^{(0)}=K\) and \(K^{(j+1)}=[K^{(j)},K^{(j)}]\), reaches the identity subgroup. For a finite soluble group \(K\), define \(F_0(K)=1\) and let \(F_{i+1}(K)\) be the inverse image of \(F(K/F_i(K))\) under the quotient map. Its Fitting height is
\[
h(K)=\min\{k\in\mathbb Z_{\geq 0}:F_k(K)=K\}.
\]
For soluble \(G\), these heights satisfy \(h^*(G)=h(G)\).

The question proposed by P. Shumyatsky is whether
\[
\forall\,G\text{ finite},\qquad
\exists\,K\leq G\quad\text{such that }K\text{ is soluble and }h(K)=h^*(G).
\]
The supplied source starts its generalized Fitting series at \(F^*_1(G)=F^*(G)\) and does not explicitly specify the height of the trivial group. The definitions above explicitly adopt the convention \(h^*(1)=h(1)=0\); with this convention the trivial-group instance is satisfied by \(K=1\).
\subsection{Short English statement}
Does every finite group have a soluble subgroup whose Fitting height equals the generalized Fitting height of the whole group? The supplied record labels this question open; that label is not independently verified here.
\subsection{Sources}
\begin{itemize}
\item Local UnsolvedMath registry, permanent ID \(701630\), exact source problem code \texttt{catalogue entry}, titled ``Kourovka Notebook Problem 19.82.'' The supplied record attributes the problem to P. Shumyatsky, proposed in 2018, and labels it open.
\item \emph{The Kourovka Notebook (Unsolved Problems in Group Theory)}, original problem 19.82, 19th issue (2018); the supplied background also identifies the 21st issue edition (2026), \texttt{arXiv:1401.0300v48}. The supplied collection description names E. I. Khukhro and V. D. Mazurov as editors. These bibliographic details are supplied attribution and have not been independently checked.
\end{itemize}
\end{document}
